\scanpage{31}\mspage{30}

For non-set-theoretic countries, one should assume the following (to be checked on the non-set-theoretic examples to come)

\margin{(?)}
% slanted
\margin{replaces \uline{Cont} 4}
\uline{Cont} $\substack{1\\4}$ For two figures \inserted{$F,\unsure{C|G}$} to be \margin{\rule{0.35pt}{4.5\baselineskip}}compatible, it is necessary and sufficient that they be so set-theoretically, and that for two \deleted{distinct strata} open strata $U$ of $F$, $V$ of $\unsure{C|G}$, $U\ne V$, \deleted{we have} and for \deleted{every} $x\in U$, $y\in V$, $\deleted{\unsure{\{x\}\{\}\{|\{x\}\}}\unsure{-|=}\{x,y\}}$ $\{\{x\},\{y\}\}$ be a figure (i.e $\{\{x\}\}$ and $\{\{y\}\}$ be compatible

% slanted
\margin{\uline{NB} The ``set-theoretic'' countries or ``countries of points'' are those which satisfy $\mathrm{Cont}'_{4}$, and such that every subset of $\mathcal L$ with 2 elements is a 0-sphere i.e. two $+$ distinct loci are disjoint there.}

\uline{NB} For every $x\in \mathcal L$, we denote by $\unsure{F|E}_x$ the figure $\{\{x\}\}$, consisting of the single stratum $\{x\}$.

% slanted
\margin{\uline{NB} A locus $x$ of $\mathcal L$ is called \uline{point-like} if if $\forall y\deleted{\ne}\in L$, $y\ne x$, \unsure{$y$} is disjoint from $x$. Thus, countries of points are those in which every locus is a point !}

\deleted{To say that the figures} If $x,y\in \mathcal L$ to say that $F_x$, $F_y$ are compatible, means that $\{\deleted{\unsure{\{x\},\{y\}}}F_x,F_y\}$ is \deleted{a} a figure. This is always true if $x=y$. \deleted{If $x\ne y$, \unsure{we then say}} \inserted{We say} that \deleted{$\{$} $x,y$ are ``\uline{disjoint}'' if they are distinct and compatible \deleted{\unsure{i.e if and only if} \unsure{$\{x$}} We also say \inserted{then} that $\{x,y\}\in \mathcal P_2(\mathcal L)$ is a 0-\uline{sphere} of $C$.

Two subsets \deleted{\unsure{of}} $A,B$ of $\mathcal L$ are called \uline{disjoint} (in the sense of the country $C$ -- or $C$-disjoint) if $\forall x\in A$, $\forall y\in B$, $x,y$ are disjoint ($\Longrightarrow A\cap B=\varnothing$).
